\documentclass[10pt,a4paper]{article}
\usepackage{amsmath,amssymb,amsthm,graphicx}
\title{Explicit asymptotically good relaxed locally correctable codes with near-log-squared query complexity}
\author{Gil Cohen, Tal Yankovitz}
\date{doi:10.4230/LIPIcs.CCC.2024.8; CC BY4.0}
\begin{document}
\maketitle
\noindent\textit{Allowed standard prerequisites (editorial).}\begin{itemize}
\item Explicit lossless conductors / unique-neighbor expander construction. Author Theorem6 citing [8,Theorem7.3]; exact quantitative statement included.
\item Standard RLCC error amplification. Author Claim4 explicitly cites [12] for its elementary proof; precise statement retained.
\end{itemize}\medskip
\section*{3 Author preliminaries}
Unless stated otherwise, all logarithms in this paper are taken to the base 2. The set of natural numbers is $\mathbb N=\{0,1,2,\ldots\}$. For $n\in\mathbb N$, $n\ge1$, we use $[n]$ to denote the set $\{1,\ldots,n\}$. For $q\in\mathbb N$, $q\ge2$, we use $H_q$ to denote the $q$-ary entropy function, and $H=H_2$ to denote the binary entropy function.
For a finite set $N$, we refer to a function $v\in\mathbb F^N$ as a vector and we say that it is indexed by $N$. For a vector $v\in\mathbb F^N$ and $i\in N$ we use $v_i$ as a shorthand for $v(i)$. For a vector $v\in\mathbb F^N$ and a set $N'\subseteq N$ we denote by $v_{N'}$ the vector $v'\in\mathbb F^{N'}$ such that $v_i'=v_i$ for every $i\in N'$. For two vectors $u,v\in\mathbb F^N$, their (absolute) hamming distance is $|\{i\in N\mid u_i\ne v_i\}|$, which we denote by $\operatorname{Dist}(u,v)$, and their relative hamming distance is $\frac{\operatorname{Dist}(u,v)}{|N|}$, which we denote by $\operatorname{RelDist}(u,v)$.
\paragraph{Definition 2.} For a finite set $N$ of size $|N|=n$ and a field $\mathbb F$, a code is a linear subspace $C\subseteq\mathbb F^N$. We say that the code $C$ is indexed by $N$ and that it is over $\mathbb F$. The length of the code is $n$. The dimension of the code, usually denoted by $k$, is the dimension of $C$ over $\mathbb F$. The (non-local) distance of the code, denoted by $d$, is $\min_{c,c'\in C,c\ne c'}\operatorname{Dist}(c,c')$. The rate of the code, typically denoted by $\rho$, is $\frac kn$. The (non-local) relative-distance of the code is defined to be $\frac dn$. The elements of $C$ are called codewords.
\paragraph{Definition 3.} A code $C\subseteq\mathbb F^N$ is called a $(q,\delta,\varepsilon)$-RLCC (relaxed locally correctable code, abbreviated) if there exists a randomized procedure $\operatorname{Cor}:\mathbb F^N\times N\to\mathbb F\cup\{\bot\}$ with the following guarantees:
\begin{itemize}
\item For every $i\in N$, $c\in C$ and $w\in\mathbb F^N$, satisfying $\operatorname{RelDist}(w,c)\le\delta$, $\operatorname{Cor}(w,i)\in\{c_i,\bot\}$ with probability at least $1-\varepsilon$.
\item $\operatorname{Cor}(c,i)=c_i$ with probability one on any $c\in C$ and $i\in N$.
\item $\operatorname{Cor}(w,i)$ always makes at most $q$ queries to $w$.
\end{itemize}
We refer to $\operatorname{Cor}$ as the local corrector (or the corrector). The parameter $\delta$ is called the correction radius, and the parameter $q$ is called the query complexity.
The error parameter of an RLCC can be easily amplified at low cost to the query complexity, as stated in the following claim (for a simple proof see, e.g., [12]).
\paragraph{Claim 4.} Let $C\subseteq\mathbb F^N$ be a $(q,\delta,\varepsilon)$-RLCC. Then, for any $h\in\mathbb N$, $C$ is also an $(hq,\delta,\varepsilon^h)$-RLCC.
We set some standard notation. Let $G=(V,E)$ be an undirected graph. For $v\in V$ we define $\Gamma(v)$ as the set of neighbors of $v$ in $G$, and let $\deg(v)$ be the degree of $v$. For a set of vertices $S\subseteq V$, we let $\Gamma(S)=\cup_{v\in S}\Gamma(v)$, and define
\[\Gamma_u(S)=\{v\in V\mid v\text{ is adjacent to exactly one }u\in S\}.\]
\paragraph{Definition 5 (Unique-neighbor expanders).} A left-$d$-regular bipartite graph $G=(L,R,E)$ is a $(\gamma,\alpha)$-unique-neighbor expander if for every $S\subseteq U$ such that $|S|\le\gamma|L|$, it holds that $|\Gamma_u(S)|\ge\alpha d|S|$.
The following theorem readily follows by the construction of lossless conductors as given by Theorem 7.3 in [8].
\paragraph{Theorem 6 ([8]).} There exist universal constants $c_0\ge1$ and $\beta\le1$ such that the following holds. For every $n$ and $m\le n$, there exists an explicit $(\gamma,\alpha)$-unique-neighbor expander $G=(L,R,E)$ with $|L|=2^n$, $|R|=2^m$, having left degree
\[d\le2^{c_0\cdot\log^3(n-m)},\]
where $\alpha=\Omega(1)$, and $\gamma=\beta\cdot\frac{2^{m-n}}d$.
\paragraph{Definition 7 (Expander codes).} Let $G=(L,R,E)$ be a bipartite graph and let $\mathbb F$ be a field. The expander code associated with $G$ is defined by
\[\operatorname{EC}_{\mathbb F}(G)=\left\{w\in\mathbb F^L\mid\forall v\in R\quad\sum_{u\in\Gamma(v)}w_u=0\right\}.\]
We usually omit the subscript $\mathbb F$ when the field is clear from context.
It is easy to see that the rate of $\operatorname{EC}_{\mathbb F}(G)$ is at least $1-\frac{|R|}{|L|}$.
\paragraph{Definition 8 (VLTCs).} A code $C\subseteq\mathbb F^N$ is called a $(q,\delta,\kappa,\sigma)$-VLTC (vicinity locally testable code, abbreviated) if there exists a randomized procedure
\[\operatorname{Tes}:\mathbb F^N\to\{\circ,\bot\}\]
with the following guarantees:
\begin{itemize}
\item For every $c\in C$ and $w\in\mathbb F^N$, satisfying $\operatorname{RelDist}(w,c)\le\delta$,
\[\Pr[\operatorname{Tes}(w)=\bot]\ge\kappa\cdot\operatorname{RelDist}(w,c)-\sigma;\]
\item $\operatorname{Tes}(c)=\circ$ with probability one on any $c\in C$.
\item $\operatorname{Tes}(w)$ always makes at most $q$ queries to $w$.
\end{itemize}
We call $\operatorname{Tes}$ a local tester (or tester for short). The parameter $q$ is referred to as the query complexity.

\paragraph{Theorem 12.} For every finite field $\mathbb F$, $n$ which is a power of 2, and $\delta>0$, there exists an explicit $(q,\delta,\frac13)$-RLCC $C\subseteq\mathbb F^{[n]}$ with query complexity
\[q=(\log n)^{2+o(1)},\]
and rate
\[\rho=1-\delta\cdot2^{O((\log\log\frac1\delta)^3)}-o(1).\]

\section*{Complete author support and main construction}
We move to show that expander codes constructed from unique-neighbor expanders are VLTCs.
\paragraph{Lemma 9.} Let $G=(L,R,E)$ be a $d$-left-regular $(\gamma,\alpha)$-unique-neighbor expander with average right-degree $\bar c$. Then, for every $b>1$, $\operatorname{EC}(G)$ is a $(b\bar c,\gamma,\alpha\bar c,\frac1b)$-VLTC.
\begin{proof}
Define
\[R'=\{v\in R\mid\deg(v)\le b\bar c\}.\]
By an averaging argument, $|R'|\ge(1-\frac1b)|R|$. The tester for $\operatorname{EC}(G)$, given oracle access to $w\in\mathbb F^L$, proceeds as follows:
\begin{enumerate}\item Sample $v\in R'$ uniformly at random.\item Query $w$ on $\Gamma(v)$.\item Output $\circ$ if $\sum_{u\in\Gamma(v)}w_u=0$; and $\bot$ otherwise.\end{enumerate}
As the sampled vertex $v$ is in $R'$, the query complexity of the tester is indeed bounded above by $b\bar c$. Further, when $w\in\operatorname{EC}(G)$, the tester outputs $\circ$ with certainty.
Consider then a word $w\in\mathbb F^L$ such that $\operatorname{RelDist}(w,c)\le\gamma$ for some codeword $c\in\operatorname{EC}(G)$. Let
\[S=\{v\in L\mid w_v\ne c_v\}.\]
As $|S|\le\gamma|L|$ we have that $|\Gamma_u(S)|\ge\alpha d|S|$. Notice that if the vertex $v$ that is sampled in Step 1 lies in $\Gamma_u(S)$ then the tester outputs $\bot$. Therefore, the probability of the tester to output $\bot$ is at least
\begin{align*}\frac{|\Gamma_u(S)|-|R\setminus R'|}{|R|}&\ge\frac{\alpha d|S|}{|R|}-\frac1b\\&=\alpha\cdot\frac{d|L|}{|R|}\cdot\frac{|S|}{|L|}-\frac1b\\&=\alpha\bar c\cdot\operatorname{RelDist}(w,c)-\frac1b,\end{align*}
which concludes the proof.
\end{proof}
We will use the following easy claim.
\paragraph{Claim 10.} Let $C\subseteq\mathbb F^N$ be a $(q,\delta,\kappa,\sigma)$-VLTC with a tester $\operatorname{Tes}$, and further let $c\in C$ and $w\in\mathbb F^N$ be such that $\alpha\le\operatorname{RelDist}(w,c)\le\delta$. Assume that we run $\operatorname{Tes}(w)$ for $g$ times, independently. Then, the probability that one of the simulations outputted $\bot$ is at least $1-e^{-\beta g}$, where $\beta=\kappa\alpha-\sigma$.
\begin{proof}
The probability that a single simulation of $\operatorname{Tes}(w)$ outputs $\bot$ is at least
\[\kappa\cdot\operatorname{RelDist}(w,c)-\sigma\ge\kappa\alpha-\sigma=\beta.\]
The probability that all the simulations output $\circ$ is thus $(1-\beta)^g\le e^{-\beta g}$.
\end{proof}
Following a similar argument to the one underlying the Kumar-Mon reduction, the following proposition states that a sequence of VLTCs can be used to construct an RLCC.
\paragraph{Proposition 11.} Let $C_1\subseteq\mathbb F^{N_1},\ldots,C_m\subseteq\mathbb F^{N_m}$ be codes with rates $\rho_1,\ldots,\rho_m$, respectively, such that for every $i\in[m-1]$, $C_i$ is a $(q',\delta',\kappa',\sigma')$-VLTC, and $C_m$ is a $(q,\delta,\kappa,\sigma)$-VLTC. Further assume that $|N_1|\le\frac1{\delta'}$, $|N_m|=n$, and for every $1<i\le m$, $|N_i|=2|N_{i-1}|$. Then, for every $g\in\mathbb N$, there exists an $((m-1)q'+gq+1,\delta,\varepsilon)$-RLCC $C\subseteq\mathbb F^{[n]}$ with rate
\[\rho\ge1-\sum_{i=1}^m(1-\rho_i),\]
where
\[\varepsilon\le1-\min\left\{\frac{\kappa'\delta'}2-\sigma',e^{g(\sigma-\frac{\kappa\delta'}2)}\right\}.\]
Moreover, if the codes $C_1,\ldots,C_m$ are explicit, then so is the resulting code $C$.
\begin{proof}
We start by describing how the code $C$ is constructed.
\paragraph{The code construction.} Let $P_1,\ldots,P_m$ be an arbitrary fixed sequence of partitions of $[n]$, satisfying that for every $i\in[m]$, $P_i$ has $2^{m-i}$ equal-size parts denoted $\{B_1^i,\ldots,B_{2^{m-i}}^i\}$, and that for every $1<i\le m$, $P_{i-1}$ is a sub-partition of $P_i$ (that is, for every $B\in P_{i-1}$, there exists $B'\in P_i$ such that $B\subseteq B'$). For every $i\in[m]$ and $B\in P_i$ let $f_{i,B}:B\to N_i$ be an arbitrary bijection, and define
\[C_{i,B}=\{c\circ f_{i,B}\mid c\in C_i\}.\]
Finally, define the code
\[C=\{w\in\mathbb F^{[n]}\mid\forall i\in[m],B\in P_i:w_B\in C_{i,B}\}.\]
The moreover part of the proof readily follows. The efficiency of the corrector will be self evident as well once the corrector is presented.
\paragraph{Rate analysis.} For every $i\in[m]$ and $B\in P_i$, the number of linear constraints required to impose so that $w_B\in C_{i,B}$ is at most $(1-\rho_i)|N_i|$. Therefore, the total number of constraints in the definition of the code $C$ is bounded above by
\[\sum_{i=1}^m|P_i|(1-\rho_i)|N_i|=\sum_{i=1}^m n(1-\rho_i),\]
which establishes the lower bound on the rate of $C$.
\paragraph{The corrector.} We turn to describe a corrector $\operatorname{Cor}:\mathbb F^{[n]}\times[n]\to\mathbb F\cup\{\bot\}$ for $C$. As for every $i\in[m]$, $C_i$ is a VLTC, it is immediate that so is $C_{i,B}$ for every $B\in P_i$, with the same parameters as $C_i$. The local tester for $C_{i,B}$ that is induced in the natural way from the local tester for $C_i$ is denoted
\[\operatorname{Tes}_{i,B}:\mathbb F^B\to\{\circ,\bot\}.\]
Let $w\in\mathbb F^{[n]}$ and $j\in[n]$. Let $r_1=r_1(j),\ldots,r_m=r_m(j)$ be the indices of blocks within the corresponding partitions $P_1,\ldots,P_m$ such that $j\in B_{r_1}^1\subseteq\cdots\subseteq B_{r_m}^m$. The corrector $\operatorname{Cor}(w,j)$ proceeds as follows:
\begin{enumerate}\item For $i=1,\ldots,m-1$, simulate $\operatorname{Tes}_{i,B_{r_i}^i}(w_{B_{r_i}^i})$.\item Simulate $\operatorname{Tes}_{m,B_{r_m}^m}(w_{B_{r_m}^m})$ for $g$ times.\item If any of the simulations outputted $\bot$, output $\bot$; otherwise, output $w_j$.\end{enumerate}
\paragraph{Query analysis.} As $\operatorname{Cor}$ simulates $m-1$ testers with query complexity $q'$, and invokes $g$ simulations of one tester with query complexity $q$, the overall query complexity is $(m-1)q'+gq+1$, accounting also for querying $w_j$.
\paragraph{Correctness.} Clearly, if $w$ is a codeword of $C$ then $w_{B_{r_i}^i}\in C_{i,B_{r_i}^i}$ for every $i\in[m]$, and so $\operatorname{Tes}_{i,B_{r_i}^i}(w_{B_{r_i}^i})=\circ$ with certainty. Therefore, $\operatorname{Cor}(w,j)=w_j$ with certainty, as required.
Assume that $w\in\mathbb F^{[n]}$ is such that $\operatorname{Dist}(w,c)\le\delta n$ for $c\in C$. Since $\operatorname{Cor}(w,j)$ always either outputs $\bot$ or $w_j$, it suffices to show that if $w_j\ne c_j$ then the corrector outputs $\bot$ with probability at least $1-\varepsilon$. Towards this end, assume $w_j\ne c_j$, and hence $w_{B_{r_1}^1}\ne c_{B_{r_1}^1}$, and further note that $w_{B_{r_m}^m}=w$ and $c_{B_{r_m}^m}=c$. For every $i\in[m-1]$ define $\delta_i=\delta'$, and let $\delta_m=\delta$. Since, per our assumption, $|N_1|\le\frac1{\delta'}$, we have that
\[\operatorname{RelDist}(w_{B_{r_1}^1},c_{B_{r_1}^1})\ge\delta'=\delta_1,\]
whereas
\[\operatorname{RelDist}(w_{B_{r_m}^m},c_{B_{r_m}^m})\le\delta=\delta_m.\]
Let $\iota\in\{2,3,\ldots,m\}$ be any index satisfying that
\begin{align*}\operatorname{RelDist}(w_{B_{r_{\iota-1}}^{\iota-1}},c_{B_{r_{\iota-1}}^{\iota-1}})&\ge\delta_{\iota-1}\\\operatorname{RelDist}(w_{B_{r_\iota}^\iota},c_{B_{r_\iota}^\iota})&\le\delta_\iota.\end{align*}
By the above account, $\iota$ is well-defined. Since $B_{r_{\iota-1}}^{\iota-1}\subseteq B_{r_\iota}^\iota$ and $|B_{r_{\iota-1}}^{\iota-1}|=\frac12|B_{r_\iota}^\iota|$, we have that
\[\frac{\delta_{\iota-1}}2\le\operatorname{RelDist}(w_{B_{r_\iota}^\iota},c_{B_{r_\iota}^\iota})\le\delta_\iota.\]
If $\iota<m$, as $\operatorname{Tes}_{\iota,B_{r_\iota}^\iota}$ is a local tester for the $(q',\delta_\iota,\kappa',\sigma')$-VLTC $C_{i,B_{r_\iota}^\iota}$ and since $c_{B_{r_\iota}^\iota}\in C_{\iota,B_{r_\iota}^\iota}$, it holds that $\operatorname{Tes}_{\iota,B_{r_\iota}^\iota}(w_{B_{r_\iota}^\iota})$ outputs $\bot$ with probability at least
\[\frac{\kappa'\delta_{\iota-1}}2-\sigma'=\frac{\kappa'\delta'}2-\sigma'.\]
If otherwise $\iota=m$ then we set
\[\beta=\frac{\kappa\delta_{m-1}}2-\sigma=\frac{\kappa\delta'}2-\sigma\]
and then by Claim 10 one of the $g$ simulations of $\operatorname{Tes}_{m,B_{r_m}^m}((w_{B_{r_m}^m}))$ with probability at least $1-e^{-\beta g}$. Thus, $\operatorname{Cor}(w,j)$ outputs $\bot$ with probability at least $\min\{\frac{\kappa'\delta'}2-\sigma',1-e^{-\beta g}\}$, as required.
\end{proof}
We are now ready to prove our main theorem.
\begin{proof}[Proof of Theorem 12]
Write $n=2^r$, and let
\[\ell=r\log r=\log n\cdot\log\log n.\]
Due to the claimed tradeoff between the rate and the correction radius, we may as well assume that $\delta\ge\frac1{2^{\lceil\log\ell\rceil}}$. We proceed to describe the sequence of expander codes that we will use, which consists of $s=r-\lceil\log\ell\rceil+1$ codes. For every $i\in[s]$, the block-length of the $i$-th code is
\[n_i=2^{\lceil\log\ell\rceil+i-1}.\]
Note that, in particular, $n_s=n$. Further, the number of linear constraints defining the $i$-th code is $m_i=2^{i-1}$ for $i\in[s-1]$, whereas the for the $s$-th code,
\[m_s=2^{r-\lfloor\log\frac1\delta+\log\beta-c_0\log^3(\log\frac1\delta)\rfloor},\]
where $\beta$ is the constant from Theorem 6.
Invoking Theorem 6, for every $i\in[s]$ let $G_i=(L_i,R_i,E_i)$ be a $d_i$-left-regular bipartite graph with $|L_i|=n_i$ and $|R_i|=m_i$ which is a $(\delta_i,\alpha)$-unique-neighbor expander for $\alpha=\Omega(1)$, such that for $i\in[s-1]$,
\begin{align*}d_i&\le2^{c_0\log^3(\lceil\log\ell\rceil)}\triangleq d,\\\delta_i&\ge\frac\beta{d\cdot2^{\lceil\log\ell\rceil}}\triangleq\delta',\end{align*}
and
\begin{align*}d_s&\le2^{c_0\log^3(\log\frac1\delta)},\\\delta_s&\ge\frac\beta{d_s\cdot2^{\lfloor\log(1/\delta)+\log(\beta)-c_0\log^3(\log(1/\delta))\rfloor}}\ge\delta.\end{align*}
The sequence of codes is defined by setting, for every $i\in[s]$, $C_i=\operatorname{EC}(G_i)$.
We turn to address the VLTC-ness of $C_1,\ldots,C_s$. Set $b=\frac4{\alpha\beta}$ and $b_s=\frac{4m_s}{\alpha\delta'd_s n_s}$. By Lemma 9, for every $i\in[s-1]$, $C_i$ is a
\[\left(\frac{bd_i n_i}{m_i}\le bd2^{\lceil\log\ell\rceil},\delta',\alpha d2^{\lceil\log\ell\rceil},\frac1b\right)\text{-VLTC},\]
and $C_s$ is a
\[\left(\frac{b_s d_s n_s}{m_s},\delta,\alpha\frac{d_s n_s}{m_s},\frac1{b_s}\right)\text{-VLTC}.\]
We further set $g=\frac{4m_s}{\alpha\delta'd_s n_s}$. We can now invoke Proposition 11 (indeed, the proposition's prerequisites are met, i.e., $n_1=2^{\lceil\log\ell\rceil}<\frac1{\delta'}$, $n_i=2n_{i-1}$, and $n_s=n$) with our choice of $g$ to obtain a code $C\subseteq\mathbb F^{[n]}$ which is an
\[\left((s-1)bd2^{\lceil\log\ell\rceil}+g\frac{b_s d_s n_s}{m_s}+1=O(sbd\ell+\delta/(\delta')^2),\delta,\varepsilon\right)\text{-RLCC},\]
where
\[\varepsilon\le1-\min\left\{\frac12\alpha d2^{\lceil\log\ell\rceil}\delta'-\frac1b,(1/e)^{g(\alpha d_s n_s\delta'/(2m_s)-1/b_s)}=(1/e)^{g\alpha d_s n_s\delta'/(4m_s)}=1/e\right\}.\]
As
\begin{align*}\frac12\alpha d2^{\lceil\log\ell\rceil}\delta'-\frac1b&=\frac{\alpha\beta}2-\frac1b\\&=\frac{\alpha\beta}2-\frac{\alpha\beta}4\\&=\frac{\alpha\beta}4,\end{align*}
we see that $\varepsilon\le1-\min\{\frac{\alpha\beta}4,1/e\}$. To decrease the error to $\frac13$, we apply Claim 4 with $h=O(1)$, and get that $C$ is also a $(q,\delta,\frac13)$-RLCC for
\[q=O(sbd\ell+\delta/(\delta')^2)=O(sd^2\ell+\delta d^2\ell^2).\]
Recall that $s\le\log n$,
\[d=2^{O((\log\log r)^3)}=2^{O((\log\log\log n)^3)}\]
and $\ell=O(\log n\cdot\log\log n)$. Therefore,
\[q=O(sd^2\ell+d^2\ell^2)=O(d^2\ell^2)=\log^2 n\cdot2^{O((\log\log\log n)^3)}=(\log n)^{2+o(1)}.\]
Lastly, as the rate $\rho_i$ of every code $C_i$ in the sequence is at least $1-\frac{m_i}{n_i}$, Proposition 11 implies that the rate $\rho$ of $C$ is lower bounded by
\begin{align*}\rho&\ge1-\sum_{i=1}^s\left(\frac{m_i}{n_i}\right)\\&=1-(s-1)\frac1{2^{\lceil\log\ell\rceil}}-\frac1{2^{\lfloor\log\frac1\delta+\log\beta-c_0\log^3(\log\frac1\delta)\rfloor}}\\&=1-O\left(\frac s\ell\right)-\delta\cdot2^{O(\log^3(\log\frac1\delta))}\\&=1-O\left(\frac r\ell\right)-\delta\cdot2^{O(\log^3(\log\frac1\delta))}\\&=1-\delta\cdot2^{O(\log^3(\log\frac1\delta))}-o(1).\end{align*}
This concludes the proof.
\end{proof}

\begin{thebibliography}{99}
\bibitem{source1} Ron Asherov and Irit Dinur. Bipartite unique-neighbour expanders via Ramanujan graphs. arXiv preprint, 2023. arXiv:2301.03072.
\bibitem{source2} Eli Ben-Sasson, Oded Goldreich, Prahladh Harsha, Madhu Sudan, and Salil Vadhan. Robust PCPs of proximity, shorter PCPs, and applications to coding. SIAM Journal on Computing, 36(4):889–974, 2006.
\bibitem{source3} Eli Ben-Sasson and Michael Viderman. Composition of semi-ltcs by two-wise tensor products. computational complexity, 24:601–643, 2015.
\bibitem{source4} Arnab Bhattacharyya, Sivakanth Gopi, and Avishay Tal. Lower bounds for 2-query LCCs over large alphabet. arXiv preprint, 2016. arXiv:1611.06980.
\bibitem{source5} Alexander R. Block, Jeremiah Blocki, Kuan Cheng, Elena Grigorescu, Xin Li, Yu Zheng, and Minshen Zhu. On relaxed locally decodable codes for Hamming and insertion-deletion errors. In 38th Computational Complexity Conference, volume 264 of LIPIcs. Leibniz Int. Proc. Inform., pages Paper No. 14, 25. Schloss Dagstuhl – Leibniz-Zentrum für Informatik, 2023. doi:10.4230/lipics.ccc.2023.14.
\bibitem{source6} Manuel Blum and Sampath Kannan. Designing programs that check their work. Journal of the ACM (JACM), 42(1):269–291, 1995.
\bibitem{source7} Clément L Canonne and Tom Gur. An adaptivity hierarchy theorem for property testing. Computational Complexity, 27(4):671–716, 2018.
\bibitem{source8} Michael Capalbo, Omer Reingold, Salil Vadhan, and Avi Wigderson. Randomness conductors and constant-degree lossless expanders. In Proceedings of the Thirty-Fourth Annual ACM Symposium on Theory of Computing, pages 659–668. ACM, New York, 2002. doi:10.1145/ 509907.510003.
\bibitem{source9} Alessandro Chiesa, Tom Gur, and Igor Shinkar. Relaxed locally correctable codes with nearly-linear block length and constant query complexity. In Proceedings of the Fourteenth Annual ACM-SIAM Symposium on Discrete Algorithms, pages 1395–1411. SIAM, 2020.
\bibitem{source10} Gil Cohen and Tal Yankovitz. Rate amplification and query-efficient distance amplification for linear LCC and LDC. In 36th Computational Complexity Conference (CCC 2021). Schloss Dagstuhl – Leibniz-Zentrum für Informatik, 2021.
\bibitem{source11} Gil Cohen and Tal Yankovitz. LCC and LDC: Tailor-made distance amplification and a refined separation. In 49th EATCS International Conference on Automata, Languages, and Programming, volume 229 of LIPIcs. Leibniz Int. Proc. Inform., pages Art. No. 44, 20. Schloss Dagstuhl – Leibniz-Zentrum für Informatik, 2022. doi:10.4230/lipics.icalp.2022.44.
\bibitem{source12} Gil Cohen and Tal Yankovitz. Relaxed locally decodable and correctable codes: beyond tensoring. In 2022 IEEE 63rd Annual Symposium on Foundations of Computer Science— FOCS 2022, pages 24–35. IEEE Computer Soc., Los Alamitos, CA, 2022.
\bibitem{source13} Itay Cohen, Roy Roth, and Amnon Ta-Shma. HDX condensers. In Electronic Colloquium on Computational Complexity (ECCC), 2023.
\bibitem{source14} Marcel Dall’Agnol, Tom Gur, and Oded Lachish. A structural theorem for local algorithms with applications to coding, testing, and privacy. In Proceedings of the 2021 ACM-SIAM Symposium on Discrete Algorithms (SODA), pages 1651–1665. SIAM, 2021.
\bibitem{source15} Irit Dinur, Shai Evra, Ron Livne, Alexander Lubotzky, and Shahar Mozes. Locally testable codes with constant rate, distance, and locality. In STOC ’22—Proceedings of the 54th Annual ACM SIGACT Symposium on Theory of Computing, pages 357–374. ACM, New York, 2022.
\bibitem{source16} Zeev Dvir, Parikshit Gopalan, and Sergey Yekhanin. Matching vector codes. SIAM Journal on Computing, 40(4):1154–1178, 2011.
\bibitem{source17} Zeev Dvir, Sivakanth Gopi, Yuzhou Gu, and Avi Wigderson. Spanoids—an abstraction of spanning structures, and a barrier for LCCs. SIAM Journal on Computing, 49(3):465–496, 2020.
\bibitem{source18} Klim Efremenko. 3-query locally decodable codes of subexponential length. SIAM Journal on Computing, 41(6):1694–1703, 2012.
\bibitem{source19} Guy Goldberg. Linear relaxed locally decodable and correctable codes do not need adaptivity and two-sided error. In Electron. Colloquium Comput. Complex., 2023.
\bibitem{source20} Oded Goldreich. Lecture notes on locally testable codes and proofs, 2016. URL: https: //www.wisdom.weizmann.ac.il/~oded/PDF/pt-ltc.pdf.
\bibitem{source21} Oded Goldreich. On the lower bound on the length of relaxed locally decodable codes. In Electronic Colloquium on Computational Complexity (ECCC), 2023.
\bibitem{source22} Oded Goldreich and Tom Gur. Universal locally verifiable codes and 3-round interactive proofs of proximity for csp. Theoretical Computer Science, 878:83–101, 2021.
\bibitem{source23} Oded Goldreich, Howard Karloff, Leonard J Schulman, and Luca Trevisan. Lower bounds for linear locally decodable codes and private information retrieval. In Proceedings 17th IEEE Annual Conference on Computational Complexity, pages 175–183. IEEE, 2002.
\bibitem{source24} Louis Golowich. New explicit constant-degree lossless expanders. arXiv preprint, 2023. arXiv:2306.07551.
\bibitem{source25} Sivakanth Gopi, Swastik Kopparty, Rafael Oliveira, Noga Ron-Zewi, and Shubhangi Saraf. Locally testable and locally correctable codes approaching the Gilbert-Varshamov bound. IEEE Transactions on Information Theory, 64(8):5813–5831, 2018.
\bibitem{source26} Alan Guo, Swastik Kopparty, and Madhu Sudan. New affine-invariant codes from lifting. In Proceedings of the 4th conference on Innovations in Theoretical Computer Science, pages 529–540. ACM, 2013.
\bibitem{source27} Tom Gur and Oded Lachish. On the power of relaxed local decoding algorithms. SIAM Journal on Computing, 50(2):788–813, 2021.
\bibitem{source28} Tom Gur, Govind Ramnarayan, and Ron Rothblum. Relaxed locally correctable codes. Theory of Computing, 16(1):1–68, 2020.
\bibitem{source29} Tom Gur and Ron D Rothblum. A hierarchy theorem for interactive proofs of proximity. In 8th Innovations in Theoretical Computer Science Conference (ITCS 2017). Schloss Dagstuhl – Leibniz-Zentrum für Informatik, 2017.
\bibitem{source30} Tom Gur and Ron D Rothblum. Non-interactive proofs of proximity. computational complexity, 27(1):99–207, 2018.
\bibitem{source31} Venkatesan Guruswami and Piotr Indyk. Linear-time encodable/decodable codes with near- optimal rate. IEEE Transactions on Information Theory, 51(10):3393–3400, 2005.
\bibitem{source32} Brett Hemenway, Rafail Ostrovsky, and Mary Wootters. Local correctability of expander codes. Information and Computation, 243:178–190, 2015.
\bibitem{source33} Jun-Ting Hsieh, Theo McKenzie, Sidhanth Mohanty, and Pedro Paredes. Explicit two-sided unique-neighbor expanders. arXiv preprint, 2023. arXiv:2302.01212.
\bibitem{source34} Jonathan Katz and Luca Trevisan. On the efficiency of local decoding procedures for error- correcting codes. In Proceedings of the thirty-second annual ACM symposium on Theory of computing, pages 80–86, 2000.
\bibitem{source35} Swastik Kopparty, Or Meir, Noga Ron-Zewi, and Shubhangi Saraf. High-rate locally correctable and locally testable codes with sub-polynomial query complexity. Journal of the ACM (JACM), 64(2):11, 2017.
\bibitem{source36} Swastik Kopparty, Shubhangi Saraf, and Sergey Yekhanin. High-rate codes with sublinear-time decoding. Journal of the ACM (JACM), 61(5):28, 2014.
\bibitem{source37} Vinayak Kumar and Geoffrey Mon. Relaxed local correctability from local testing. In Electron. Colloquium Comput. Complex., 2023.
\bibitem{source38} Richard J Lipton. Efficient checking of computations. In Annual Symposium on Theoretical Aspects of Computer Science, pages 207–215. Springer, 1990.
\bibitem{source39} Or Meir. Locally correctable and testable codes approaching the Singleton bound. In Electronic Colloquium on Computational Complexity (ECCC), volume 21(107), page 14, 2014.
\bibitem{source40} Dana Moshkovitz and Ran Raz. Two-query PCP with subconstant error. Journal of the ACM (JACM), 57(5):1–29, 2008.
\bibitem{source41} Pavel Panteleev and Gleb Kalachev. Asymptotically good quantum and locally testable classical LDPC codes. In STOC ’22—Proceedings of the 54th Annual ACM SIGACT Symposium on Theory of Computing, pages 375–388. ACM, New York, 2022.
\bibitem{source42} Noga Ron-Zewi and Ron D Rothblum. Local proofs approaching the witness length. In 2020 IEEE 61st Annual Symposium on Foundations of Computer Science (FOCS), pages 846–857. IEEE, 2020.
\bibitem{source43} Michael Sipser and Daniel A. Spielman. Expander codes. IEEE Trans. Inform. Theory, 42(6, part 1):1710–1722, 1996. Codes and complexity. doi:10.1109/18.556667.
\bibitem{source44} Daniel Alan Spielman. Computationally efficient error-correcting codes and holographic proofs. ProQuest LLC, Ann Arbor, MI, 1995. Thesis (Ph.D.) – Massachusetts Institute of Technology. URL: http://gateway.proquest.com/openurl?url\_ver=Z39.88-2004\&rft\_val\_fmt=info: ofi/fmt:kev:mtx:dissertation\&res\_dat=xri:pqdiss\&rft\_dat=xri:pqdiss:0576626.
\bibitem{source45} Sergey Yekhanin. Towards 3-query locally decodable codes of subexponential length. Journal of the ACM (JACM), 55(1):1–16, 2008.
\end{thebibliography}
\end{document}
